% ------------------------
% ---- DOCUMENT STYLE ----
% ------------------------

\pagestyle{plain}
\raggedbottom % Se la pagina non è completa, lascia lo spazio alla fine
\titleformat{\chapter}[display]
    {\normalfont\huge\bfseries}{\chaptertitlename\ \thechapter}{10pt}{\LARGE}
\titlespacing*{\chapter}{0pt}{0pt}{20pt}
\chaptertitlefont{\fontsize{22pt}{30pt}\selectfont}

%\hypersetup{ % Setup dell'aspetto dei link, se vuoi qualcosa di particolare
%    colorlinks,
%    citecolor=black,
%    filecolor=black,
%    linkcolor=black,
%    urlcolor=black
%}

% \renewcommand{\footnoterule}{ % Rende la linea delle footnotes larga tutta la pagina
%   \kern -3pt
%   \hrule width \textwidth height 1pt
%   \kern 2pt
% } 

\renewcommand{\footnotesize}{\fontsize{11pt}{13pt}\selectfont} % Imposta la dimensione del testo delle footnotes
\setlength{\footnotesep}{0.5cm} % Imposta lo spazio fra e singole footnotes
\setlength{\skip\footins}{1.5cm} % Imposta lo spazio fra il corpo e le footnotes

% Definizione colori
\definecolor{codebackground}{gray}{0.96}      
\definecolor{commentgrey}{HTML}{888888}      
\definecolor{keywordblue}{HTML}{1C82B6}     
\definecolor{funcred}{HTML}{D3545F}        
\definecolor{stringgreen}{HTML}{7CA736}          
\definecolor{returnorange}{HTML}{CC5500}

% Configurazione globale
\lstset{
    language=C++,
    basicstyle=\ttfamily\scriptsize,          
    backgroundcolor=\color{codebackground},
    breaklines=true,
    breakatwhitespace=true,
    showstringspaces=false,
    keepspaces=true,
    frame=single,
    framerule=0pt,
    stringstyle=\color{stringgreen},
    commentstyle=\color{commentgrey},
    alsoletter={_},
    keywordstyle=\color{keywordblue},
    morekeywords={using, for, if, else, int, float, double, char, void, bool, struct, class, public, private, protected, const, static, virtual, template, typename, switch, case, default, break, continue, do, while, new, delete, try, catch, throw, inline, typedef, sizeof, enum, extern, volatile, register, union, explicit, friend, operator, this, goto, nullptr, auto, constexpr, decltype, true, false, namespace},
    emph={main, printf, MPI_Init, MPI_Finalize, should_exploit_symmetry, rows_to_compute, escape_time, mirror_computed_rows, lies_in_main_cardioid, lies_in_period_two_bulb, pruning_enabled, symmetry_enabled, compute_mandelbrot, MandelbrotImage, Viewport},
    emphstyle=\color{funcred},
    emph={[2]return},
    emphstyle={[2]\color{returnorange}},
    % Mappatura direttive preprocessore e header
    literate={\#include}{{\textcolor{keywordblue}{\#include}}}{8}
             {\#define}{{\textcolor{keywordblue}{\#define}}}{7}
             {<cmath>}{{\textcolor{stringgreen}{<cmath>}}}{7} 
             {<cstddef>}{{\textcolor{stringgreen}{<cstddef>}}}{9} 
             {<stdid.h>}{{\textcolor{stringgreen}{<stdid.h>}}}{9} 
             {<stdio.h>}{{\textcolor{stringgreen}{<stdio.h>}}}{9}
}

\lstdefinelanguage{CUDA}{
    language=C++,
    morekeywords={
        % execution configuration
        __global__, __device__, __host__, __shared__, __constant__,
        % built-in variables
        threadIdx, blockIdx, blockDim, gridDim, warpSize,
        % types
        dim3, cudaError_t, cudaStream_t, cudaEvent_t,
        % memory functions
        cudaMalloc, cudaFree, cudaMemcpy, cudaMemset,
        cudaMallocManaged, cudaDeviceSynchronize,
        cudaMemcpyHostToDevice, cudaMemcpyDeviceToHost,
        % occupancy API
        cudaOccupancyMaxActiveBlocksPerMultiprocessor,
        cudaOccupancyMaxPotentialBlockSize,
        % error handling
        cudaGetLastError, cudaGetErrorString, cudaSuccess
    },
    sensitive=true
}

% Definizione ambiente inline
\newtcbox{\inlinecode}[1][codebackground]{
    on line,
    box align=base,
    colback=#1,
    colframe=#1,
    boxrule=0pt,
    arc=3pt,
    outer arc=3pt,
    leftrule=0pt,
    rightrule=0pt,
    toprule=0pt,
    bottomrule=0pt,
    boxsep=0pt,
    left=2pt,
    right=2pt,
    top=2pt,
    bottom=2pt
}

\newcommand{\code}[1]{\inlinecode{\lstinline{#1}}}

\captionsetup{labelfont=bf}

